%%%%%%%%%%%%%%%%%%%%%%%%%%%%%%%%%%%%%%%%%%%%%%%%%%%%%%%%%%%%%%%%%%%%%%
% Lớp 10 | Thống kê | Các số đặc trưng đo xu thế trung tâm của mẫu số liệu
% Dạng 0D6?3-4: Tứ phân vị 
% Số câu: 6  (sinh tự động bởi tach_cau.py; từ nay thầy sửa trực tiếp file này)
%%%%%%%%%%%%%%%%%%%%%%%%%%%%%%%%%%%%%%%%%%%%%%%%%%%%%%%%%%%%%%%%%%%%%%

%%%%%%%%%%%%%%%% TRẮC NGHIỆM 4 PHƯƠNG ÁN %%%%%%%%%%%%%%%%

% nguồn: MaHoaNH_TracNghiem_Lop10
\begin{ex}%[0D6N3-4]%[ma:eac49cbb]
	Số học sinh của các lớp khối $10$ trong một trường trung học được thống kê như sau:\begin{center}
	\begin{tabular}{cccccccc}
		35 & 40 & 37 & 46 & 38 & 42 &  43 & 36
	\end{tabular}
\end{center} Tứ phân vị thứ hai của mẫu số liệu được thống kê ở trên là
\choice 
	{$40$} 
	 {\True $39$}
	{$36$} 
	 {$42$}
\loigiai{ 
	Sắp xếp số liệu theo thứ tự không giảm: \begin{center}
	\begin{tabular}{cccccccc}
		35 & 36 &  37 & 38 & 40 & 42 & 43 & 46.
	\end{tabular}
	\end{center}
Dãy số liệu trên có $2$ số chính giữa nên  tứ phân vị thứ hai là $Q_2=\dfrac{38+40}{2}=39$.}
\end{ex}

% nguồn: MaHoaNH_TracNghiem_Lop10
\begin{ex}%[0D6H3-4]%[ma:d286b7be]
	Tiền lương hàng tháng của $ 7 $ nhân viên trong một công ty du lịch lần lượt là  $ 6{,}0 $; $ 8{,}5 $; $ 7{,}0 $; $ 7{,}5 $; $ 2{,}5 $; $ 6{,}7 $; $ 3{,}0 $ (đơn vị triệu
	đồng). Khoảng tứ phân  vị của dãy số liệu thống kê trên bằng
	\choice 
		{$ 3{,}0 $ triệu đồng} 
		 {$ 6 $ triệu đồng}
		{\True $  4{,}5 $ triệu đồng}
		{$ 6{,}7 $ triệu đồng} 
	 \loigiai{
		Sắp xếp các số liệu theo tăng dần, ta được 
		\begin{center}
			$  2{,}5 $; $ 3{,}0$; $6{,}0 $; $ 6{,}7 $; $ 7{,}0 $; $ 7{,}5 $; $ 8{,}5 $
		\end{center}
		Do đó $ \mathrm{Q}_1=3{,}0 $; $  \mathrm{Q}_2=6{,}7 $; $ \mathrm{Q}_3=7{,}5 $.\\
		Vậy khoảng tứ phân vị của mẫu số  liệu trên là $ \Delta_{\mathrm{Q}}=\mathrm{Q}_3-\mathrm{Q}_1= 7{,}5-3{,}0=4{,}5$ (triệu đồng).
	 }
\end{ex}

% nguồn: MaHoaNH_TracNghiem_Lop10
\begin{ex}%[0D6H3-4]%[ma:84acd4a0]
	Tìm tứ phân vị thứ ba của mẫu số liệu: \begin{center}
	\begin{tabular}{ccccccccccc}
		27 & 15 & 18 & 30 & 19 &  40 & 100 & 9 & 46 & 10 & 200.
	\end{tabular}
\end{center}
\choice 
	{$18$} 
	{$15$} 
	{$40$} 
	 {\True $46$}
 \loigiai{
	Sắp xếp số liệu theo thứ tự không giảm: \begin{center}
	\begin{tabular}{ccccccccccc}
		9 & 10 & 15 & 18 & 19 & 27 & 30 & 40 &  46 & 100 & 200.
	\end{tabular}
	\end{center}
	Tứ phân vị thứ nhì là trung vị của  dãy số liệu là $Q_2=27$.\\
	Tứ phân vị thứ ba là trung vị của dãy số liệu: \begin{center}
	\begin{tabular}{ccccc}
		30 & 40 & 46 & 100 &  200.
	\end{tabular}
	\end{center}
Khi đó tứ phân vị thứ  ba là $Q_3=46$.}
\end{ex}


%%%%%%%%%%%%%%%% TRẢ LỜI NGẮN %%%%%%%%%%%%%%%%

% nguồn: MaHoaNH_TraLoiNgan_Lop10
\begin{bt}%[0D6H3-4]%[ma:313e58b4]
	Hãy tìm các tứ phân vị của mẫu số liệu sau $2$; $3$; $7$; $4$; $5$; $3$; $4$; $4$. Khi đó $Q_1+Q_2+Q_3$ có giá  trị bằng bao nhiêu?
	\shortans{11,5} 
	\loigiai{ 
		Sắp xếp mẫu $2$; $3$; $3$; $4$; $3$; $3$; $5$; $7$. Kích thước mẫu là $8$ (chẵn).\\ 
		Khi  đó trung vị của mẫu là trung bình cộng của giá trị thứ $4$ và $5$ là $M_e=\dfrac{1}{2}(4+4)=4$.\\
		Tứ phân vị thứ nhất là  $Q_1=\dfrac{1}{2}(3+3)=3$.\\
		Tứ phân  vị thứ hai là $Q_2=M_e=4$.\\
		Tứ  phân vị thứ ba là $Q_3=\dfrac{1}{2}(4+5)=4{,}5$.\\
		 Vậy $Q_1+Q_2+Q_3=11{,}5$.
	} 
\end{bt}


%%%%%%%%%%%%%%%% TỰ LUẬN %%%%%%%%%%%%%%%%

% nguồn: MaHoaNH_TuLuan_Lop10
\begin{bt}%[0D6H3-4]%[ma:ba5ffd9a]
	Bảng sau đây cho biết số lần  học tiếng Anh trên internet trong một tuần của một học sinh lớp $10$
	\begin{center}
		\begin{longtable}{|l|l|l|l|l|l|l|}
			 \hline
			Số lần & $0$  & $1$ & $2$ & $3$ & $4$ & $5$\\
			\hline 
			Số học  sinh & $2$ & $4$ & $6$ & $12$ & $8$ & $3$\\
			\hline 
		\end{longtable}
	\end{center}
	Hãy tìm các tứ phân vị cho  mẫu số liệu này. Tính $T=Q_{1}+Q_{2}+Q_{3}$.
	\loigiai{ 
		Cỡ mẫu  $n=2+4+6+12+8+3=35$.\\
		Vì cỡ mẫu là $n=35$ là số lẻ nên giá trị của tứ phân vị thứ hai  là $Q_2=x_{18}=3$.\\
		Tứ phân vị thứ nhất là trung vị của mẫu: $x_1; x_2; \ldots; x_{17}$. Do đó $Q_1=x_9=2$.\\ 
		Tứ phân vị thứ  ba là trung vị của mẫu: $x_{19}; x_{20}; \ldots; x_{35}$. Do đó $Q_3=x_{27}=4$.\\
		 Vậy $T=Q_{1}+Q_{2}+Q_{3}=9$.
	 }
\end{bt}

% nguồn: MaHoaNH_TuLuan_Lop10
\begin{bt}%[0D6H3-4]%[ma:34ad2959]
	Kết  quả bài kiểm tra môn toán của một nhóm học sinh được ghi lại như sau
	 $$
	\begin{array}{lllllllllllll}
		8 & 6 & 9 & 10 & 5 & 7 & 9 & 8 & 7 & 6 & 8  & 8 & 10
	\end{array}
	$$ 
	Tìm các tứ phân vị của  mẫu số liệu trên.
	 \loigiai{
		Xếp mẫu số liệu theo thứ  tự không giảm, ta được
		 $$
		\begin{array}{lllllllllllll}
			5 & 6 & 6 & 7  & 7 & 8 & 8 & 8 & 8 & 9 & 9 & 10 & 10
		\end{array}
		$$ 
		Mẫu số liệu  có $13$ giá trị, do đó
		\begin{itemize}
			 \item $Q_2=x_{7}=8$;
			\item $Q_1=\dfrac{1}{2}\left(x_3+x_4\right)=\dfrac{1}{2}\cdot  (6+7)=\dfrac{13}{2}$;
			 \item $Q_3=\dfrac{1}{2}\left(x_{10}+x_{11}\right)=\dfrac{1}{2}\cdot (9+9)=9$.
		\end{itemize}
	} 
\end{bt}
